% !TeX root = ../Cosmology.tex

%% --------------------------------------------------------
\chapter{Гідродинаміка фотон-баріонної плазми\\ та акустичні осциляції}
\label{sec:bao}
%% --------------------------------------------------------

Після розігріву Всесвіт входить у радіаційно-доміновану епоху.
Температура настільки висока, що баріонна матерія повністю
іонізована: протони та електрони утворюють плазму, в якій
фотони інтенсивно розсіюються на вільних електронах
(томсонівське розсіяння). Середня довжина вільного пробігу
фотона набагато менша за космологічний горизонт, тому
випромінювання і баріони поводяться як єдина
\emph{фотон-баріонна плазма} \cite{mukhanov_2005}.

Первинні неоднорідності кривини $\mathcal{R}_k$, успадковані
від інфляції, стають початковими умовами для цієї плазми.
Подальша динаміка визначається конкуренцією двох сил:
гравітаційне притягання намагається стиснути речовину у
потенціальні ями $\Phi$, а внутрішній тиск фотонного газу
чинить опір. Результат~--- акустичні коливання.

Перш ніж переходити до суворих кількісних рівнянь релятивістської
кінетики, уявимо фізичну картину цих процесів наочно.

Космологічне середовище того часу можна порівняти з гігантським
пружним батутом. Роль самого полотна батута відіграє тканина
простору-часу (гравітаційне поле). Первинні збурення кривини
$\mathcal{R}_k$, які Всесвіт успадкував від інфляції~--- це початкові
нерівності, «вм'ятини» та «горби» на самому полотні цього батута.

Холодна темна матерія, яка майже не має власного тиску і є абсолютно
«прозорою» для фотонного випромінювання, намагається реагувати на
первинні нерівності $\mathcal{R}_k$ першою.
Однак у радіаційно-доміновану епоху зростанню її збурень заважає
так звана \emph{стагнація Месароша} \cite{Meszaros1974}: оскільки
густина енергії випромінювання переважає, параметр Хаббла
визначається ним, і це гальмує гравітаційний колапс темної
матерії~--- збурення її густини ростуть лише логарифмічно по $a$,
а не степенево.
Ситуація корінним чином змінюється після епохи рівності матерії та
випромінювання ($z_{\text{eq}} \approx 3400$). Скинувши ці обмеження,
темна матерія починає домінувати гравітаційно. Вона безперешкодно
злітається на дно вм'ятин і своєю масою остаточно продавлює
«батут» простору-часу, перетворюючи первинний рельєф на стабільні,
нерухомі «гравітаційні колодязі»~--- потенціали $\Phi$.

Потрапляючи в ці вже підготовлені структури, фотон-баріонна плазма,
що заперта власним тиском, поводить себе як динамічна система з двох
міцно склеєних компонентів:
\begin{itemize}
    \item \textbf{Фотони} працюють як суперпружна, невагома \emph{пружина}.
    Вони створюють колосальний тиск, чинять опір будь-якому стисненню
    і прагнуть розштовхати середовище назад.
    \item \textbf{Баріони} (протони та електрони) працюють як важкі
    \emph{цеглини} або мішки з піском. Самі по собі вони не мають
    пружності, але забезпечують системі масу та інерцію.
\end{itemize}
Коли такий згусток плазми опиняється в гравітаційній ямі, зовнішнє
поле $\Phi$ затягує важкі баріони до центру (фаза стиснення).
Проте, стискаючись, фотонна «пружина» накопичує тиск і в певний
момент із силою виштовхує речовину назад (фаза розширення).
Цей циклічний процес акустичних коливань триватиме всередині
готових колодязів аж до моменту рекомбінації.


Тепер, тримаючи в голові цю механічну аналогію, математично опишемо,
як саме виводяться ці закони руху із тензора енергії-імпульсу.

\medskip
\noindent\textbf{Логічна карта розділу:}
\[
\mathcal{R}_k \;\xrightarrow[\text{надгоризонт.}]{\text{радіац. епоха}}\;
\Phi_k \;\xrightarrow[\text{вхід під горизонт}]{\text{рушійна сила}}\;
\Theta_0(k,\eta) \;\xrightarrow[\eta=\eta_{\text{rec}}]{\text{«стоп-кадр»}}\;
\mathcal{D}_\ell^{TT}\ \&\ \mathrm{BAO}.
\]

\medskip
\noindent\textbf{Крок 1: $\mathcal{R}_k \to \Phi_k$~--- спадщина інфляції.}
Первинна кривина $\mathcal{R}_k$, заморожена під час інфляції,
після реогріву перетворюється на класичний ньютонівський
гравітаційний потенціал. \textbf{У радіаційно-доміновану епоху,
на надхабблових масштабах} ($k \ll aH$) співвідношення має вигляд:
\[
\Phi_k = -\frac{2}{3}\,\mathcal{R}_k.
\]
Це «гравітаційні ями» і «горби», успадковані від квантового
шуму~--- початкові умови для всієї подальшої динаміки.

\medskip
\noindent\textbf{Крок 2: $\Phi_k \to \Theta_0(k,\eta)$~--- акустичний осцилятор.}
Фотон-баріонна плазма реагує на ці ями конкуренцією двох сил:
гравітаційне притягання стискає речовину, фотонний тиск~---
розштовхує. Результат~--- гармонічний осцилятор із рівнянням
\[
\Theta_0'' + c_s^2 k^2\Theta_0 = -\frac{k^2}{3}\Phi,
\qquad
c_s = \frac{1}{\sqrt{3(1+R)}},
\]
де $\Theta_0 = \delta T/T$~--- температурна флуктуація, а
$c_s$~--- швидкість звуку у плазмі.
Кожна мода $k$ коливається незалежно, але всі стартують
з однакової фази косинуса~--- наслідок когерентності
інфляційних початкових умов.

\medskip
\noindent\textbf{Крок 3: $\Theta_0(k,\eta_{\text{rec}})$~--- «стоп-кадр».}
У момент рекомбінації ($z \approx 1100$) фотони відриваються
від баріонів і вільно летять до нас. Акустичні коливання
«заморожуються» у тому стані, в якому їх застала рекомбінація.
Максимальна відстань, яку встигла пройти звукова хвиля~---
звуковий горизонт $r_s \approx 147$~Мпк~--- є єдиним
лінійним масштабом у задачі.

\medskip
\noindent\textbf{Крок 4: $\Theta_0(k,\eta_{\text{rec}}) \to \mathcal{D}_\ell^{TT}$~--- проекція на небо.}
Просторове поле $\Theta_0(\mathbf{r},\eta_{\text{rec}})$ проектується
на небесну сферу: моди, що в момент рекомбінації були у точці
максимального стиснення або розширення, дають піки кутового
спектра $\mathcal{D}_\ell^{TT}$. Положення піків визначається
виключно $r_s$, а їхні відносні висоти~--- баріонною густиною
$\Omega_b h^2$ і темною матерією $\Omega_c h^2$.

\medskip
\noindent\textbf{Крок 5: BAO~--- той самий $r_s$ у галактиках.}
Той самий звуковий горизонт залишає слід і в баріонній
речовині: після рекомбінації баріони «замерзають» на
радіусі $r_s$ навколо кожної первинної неоднорідності.
Сьогодні це проявляється як надлишок у кореляційній
функції галактик на відстані $\sim 147$~Мпк~---
«стандартна лінійка» для вимірювання $H(z)$.


%% --------------------------------------------------------
\section{Рівняння руху матерії у збуреній метриці}
\label{sec:matter_perturbations}
%% --------------------------------------------------------

Динаміка будь-якої речовини у викривленому просторі-часі
визначається законом збереження тензора енергії-імпульсу:
\begin{equation}\label{eq:conservation_covariant}
  \nabla_\mu T^{\mu\nu} = 0.
\end{equation}
Для ідеальної рідини з густиною $\rho$, тиском $p$ і
4-швидкістю $u^\mu$:
\begin{equation}\label{eq:Tmunu_fluid}
  T^{\mu\nu} = (\rho + p)\,u^\mu u^\nu + p\,g^{\mu\nu}.
\end{equation}
Розкладемо кожну величину на однорідний фон і мале збурення:
\begin{equation*}
  \rho = \bar\rho + \delta\rho, \qquad
  p    = \bar p   + \delta p,   \qquad
  u^\mu = \bar u^\mu + \delta u^\mu,
\end{equation*}
де фонова 4-швидкість для співрухомого спостерігача
$\bar u^\mu = (1/a,\, \mathbf{0})$ у конформному часі.
Збурення просторової швидкості позначимо
$\delta u^i \equiv v^i/a$\footnote{%
Множник $1/a$ має простий геометричний зміст.
Просторові координати $x^i$ у метриці FLRW є \emph{супутніми}~---
вони не змінюються при хаббловському розширенні.
Фізична пекулярна швидкість (відхилення від хаббловського потоку)
дорівнює $v^i = a\dot{x}^i$, тобто $\dot{x}^i = v^i/a$.
Оскільки просторова компонента 4-швидкості у конформному часі
$\delta u^i = dx^i/d\eta = a\,\dot{x}^i/a = \dot{x}^i$~---
виходить $\delta u^i = v^i/a$.
Інакше кажучи: $v^i$~--- це фізична пекулярна швидкість,
а $\delta u^i$~--- її аналог у конформних координатах,
послаблений множником $a$ через розтягування просторової сітки.%
}, де $\mathbf{v}$~--- фізична
пекулярна швидкість рідини.

%% --------------------------------------------------------
\subsection*{Рівняння неперервності ($\nu = 0$)}
%% --------------------------------------------------------

Нульова компонента~\eqref{eq:conservation_covariant}
у конформно-ньютонівській метриці~\eqref{eq:perturbed_metric}
після лінеаризації дає:
\begin{equation}\label{eq:continuity_general}
  \delta' + (1+w)\left(\theta - 3\Psi'\right)
  + 3\mathcal{H}\left(c_s^2 - w\right)\delta = 0,
\end{equation}
де $\delta \equiv \delta\rho/\bar\rho$~--- відносний контраст
густини, $\theta \equiv \partial_i v^i$~--- дивергенція поля
швидкостей, $w \equiv \bar p/\bar\rho$~--- параметр рівняння
стану, $c_s^2 \equiv \delta p/\delta\rho$~--- квадрат швидкості
звуку\footnote{%
Для адіабатичних збурень (які ми розглядаємо) $c_s^2 = \dot{\bar p}/\dot{\bar\rho}$;
для фотон-баріонної плазми ця умова виконується.}.
Член $(1+w)(\theta - 3\Psi')$~--- потік речовини з
урахуванням метричного збурення; $3\mathcal{H}(c_s^2 - w)\delta$~---
хаббловське розбавлення густини.

Для фотонів ($w = 1/3$, $c_s^2 = 1/3$,
$\delta_\gamma \equiv \delta\rho_\gamma/\bar\rho_\gamma$)
і на субхабблових масштабах де $\Psi' \ll k^2\Psi$,
рівняння~\eqref{eq:continuity_general} у Фур'є-просторі
спрощується до:
\begin{equation}\label{eq:continuity_photon}
  \delta_\gamma' = -\frac{4}{3}\,\theta + 4\Phi',
\end{equation}
де член $4\Phi'$ відображає зміну об'єму просторових
гіперповерхонь через потенціал $\Phi$.

%% --------------------------------------------------------
\subsection*{Рівняння Ейлера ($\nu = i$)}
%% --------------------------------------------------------

Просторові компоненти~\eqref{eq:conservation_covariant}
після лінеаризації дають рівняння руху для поля швидкостей.
Фізичний зміст кожного члена стає прозорим, якщо записати
результат у вигляді узагальненого другого закону Ньютона для
елемента рідини:
\begin{equation}\label{eq:euler_general}
  \theta' + \mathcal{H}(1 - 3w)\,\theta
  = -\frac{c_s^2}{1+w}\,k^2\delta - k^2\Phi.
\end{equation}
Тут:
\begin{itemize}
  \item $\theta'$~--- прискорення елемента рідини,
  \item $\mathcal{H}(1-3w)\,\theta$~--- хаббловське
        гальмування (тертя через розширення простору):
        для нерелятивістської матерії ($w=0$) воно сповільнює
        рух, для радіації ($w=1/3$) воно зникає,
  \item $-\dfrac{c_s^2}{1+w}\,k^2\delta$~---
        сила градієнта тиску: саме вона є \emph{відновлювальною
        силою} акустичних коливань,
  \item $-k^2\Phi$~--- гравітаційне притягання до потенціальних
        ям~--- сила що намагається стиснути речовину.
\end{itemize}

Для об'єднаної фотон-баріонної рідини з параметром
$R \equiv 3\bar\rho_b/4\bar\rho_\gamma$\footnote{%
Множник $3/4$ не є довільним: у релятивістській
гідродинаміці роль інертної маси відіграє ентальпія
$\rho + p$, а не лише $\rho$. Для фотонного газу
$p_\gamma = \rho_\gamma/3$, тому
$\rho_\gamma + p_\gamma = \tfrac{4}{3}\rho_\gamma$,
і параметр $R$ є відношенням інерційної маси баріонів
до інерційної маси фотонів:
$R = \bar\rho_b / (\tfrac{4}{3}\bar\rho_\gamma)
   = 3\bar\rho_b / 4\bar\rho_\gamma$.
Власна швидкість звуку баріонного газу при
$T \sim 10^4$~К нерелятивістська:
$c_s^{(b)} \sim \sqrt{T/m_p} \sim 10^{-5}\,c$~---
тому баріонний тиск нехтовний і тиск плазми
визначається виключно фотонами.%
}
рівняння Ейлера~\eqref{eq:euler_general} після лінеаризації
та врахування тісного зчеплення набуває вигляду:
\begin{equation}\label{eq:euler_photon_baryon}
  \theta' + \mathcal{H}\frac{R}{1+R}\,\theta
  = \frac{k^2}{1+R}\,\Theta_0 - k^2\Phi,
\end{equation}
де $\Theta_0 \equiv \delta_\gamma/4 = \delta T/T$~---
температурна флуктуація.

Фізичний зміст окремих членів:
\begin{itemize}
  \item Член $\mathcal{H}\frac{R}{1+R}\,\theta$ описує
        ефективне тертя, зумовлене розширенням простору.
        Він менший, ніж у чистій фотонній плазмі,
        оскільки баріони додають інерцію, але не
        збільшують тертя пропорційно.
  \item Член $\frac{k^2}{1+R}\,\Theta_0$ походить від
        градієнта фотонного тиску. Велика інерція
        баріонів ($R>0$) зменшує прискорення, яке
        створює цей градієнт~--- тому в знаменнику
        з'являється $1+R$.
  \item Член $-k^2\Phi$~--- звичайне гравітаційне
        притягання до потенціальної ями.
\end{itemize}

%% --------------------------------------------------------
\subsection*{Від двох рівнянь до акустичного осцилятора}
%% --------------------------------------------------------

Акустичні коливання виникають лише на \emph{субхабблових}
масштабах ($k \gg \mathcal{H}$): тільки тоді різні точки
плазми перебувають у причинному контакті і можуть
коливатися узгоджено. На надхабблових масштабах
($k \ll \mathcal{H}$) моди заморожені.

Щоб побачити структуру осцилятора явно, виключимо
$\theta$ з системи~\eqref{eq:continuity_photon}
і~\eqref{eq:euler_photon_baryon}: диференціюємо перше
рівняння за $\eta$ і підставляємо $\theta'$ з другого.
У конформно-ньютонівському калібруванні $\Phi = \Psi$;
нехтуємо членами з $\mathcal{H}$ на субхабблових
масштабах~--- отримуємо рівняння вимушеного акустичного
осцилятора:
\begin{equation}\label{eq:acoustic_preview}
  \frac{d}{d\eta}\bigl[(1+R)\,\delta_\gamma'\bigr]
  + \frac{k^2}{3(1+R)}\,\delta_\gamma
  = -\frac{k^2}{3}\,\Phi.
\end{equation}
\textbf{Ліва частина} --- осцилятор із змінною ефективною масою
$(1+R)$ у першому члені (похідна) та частотою, що залежить
від $1/(1+R)$.
Чому саме $(1+R)$? У релятивістській рідині
відновлювальна сила визначається тиском, а інерція~---
ентальпією $\rho + p$, а не лише $\rho$.
Для фотон-баріонної суміші ефективна інерція є
$\rho_\gamma + p_\gamma + \rho_b = \tfrac{4}{3}\rho_\gamma(1+R)$,
звідси швидкість звуку:
\begin{equation}\label{eq:sound_speed}
  c_s(\eta) = \frac{1}{\sqrt{3\bigl(1+R(\eta)\bigr)}}.
\end{equation}
При $R \to 0$: $c_s \to 1/\sqrt{3}$~--- чиста фотонна
плазма. При $R \gg 1$: $c_s \to 0$~--- важкі баріони
гальмують коливання. Оскільки $R \propto a$, швидкість
звуку монотонно спадає до рекомбінації.

\textbf{Права частина}~--- гравітаційне вимушення
потенціалом $\Phi$, що зміщує рівноважне положення
осцилятора: баріони «скочуються» у ями темної матерії,
що підсилює стиснення і підвищує непарні піки відносно
парних.

Зауважимо, що рівняння~\eqref{eq:acoustic_preview}
записане лише для фотонної компоненти $\delta_\gamma$~---
але це не спрощення.
У режимі тісного зчеплення~--- коли фотони встигають
розсіятися на електронах багато разів за один
хаббловський час ($\tau_c \ll H^{-1}$)~--- томсонівське розсіяння утримує
фотони і баріони як єдину рідину:
$\delta_b \approx \tfrac{3}{4}\delta_\gamma$  і
поля пекулярних швидкостей двох компонент збігаються
($\theta_b = \theta_\gamma$, де $\theta \equiv \partial_i v^i$~---
дивергенція поля швидкостей). Вплив баріонів
повністю закодований у параметрі $R$~--- через
ефективну інерцію $(1+R)$ і швидкість звуку
$c_s = 1/\sqrt{3(1+R)}$.

Кожна мода $k$ коливається незалежно. Початкова умова
для всіх мод задається інфляційним збуренням $\mathcal{R}_k$  (§\,\ref{sec:quantum}):
при вході під горизонт $\delta_\gamma \propto \mathcal{R}_k$
і $\delta_\gamma' \approx 0$~--- тобто всі моди стартують
з максимального стиснення, у фазі косинуса.
Розв'язок осцилятора тому має вигляд
$\delta_\gamma(k,\eta) \propto \mathcal{R}_k\cos(k r_s(\eta))$,
де $r_s = \int_0^\eta c_s\,d\eta'$~--- звуковий горизонт.
Оскільки початкова фаза однакова для всіх мод з однаковим $k$
незалежно від їхнього просторового положення, осциляції є
\emph{когерентними}~--- саме це породжує чіткі акустичні
піки на спектрі CMB, на відміну від некогерентних джерел
типу топологічних дефектів, які дають лише згладжений горб
(рис.~\ref{fig:acoustic_modes_A}).

%% ======================================================================
%% Рисунок Б: просторовий знімок при η_rec
%% ======================================================================
\begin{figure}[h!]
\centering
\definecolor{modeone}{RGB}{232,89,60}
\definecolor{modetwo}{RGB}{29,158,117}
\definecolor{modethree}{RGB}{127,119,221}
\definecolor{modesum}{RGB}{186,117,23}
\definecolor{recline}{RGB}{31,119,180}

\begin{tikzpicture}
\begin{axis}[
    width  = \linewidth,
    height = 0.5\linewidth,
    xmin = 0, xmax = 580,
    ymin = -1.5, ymax = 1.5,
    xlabel = {Просторова координата $r$ [Мпк]},
    ylabel = {$\delta_\gamma(r)$},
    label style     = {font=\small},
    tick label style= {font=\footnotesize},
    axis lines = left,
    xtick = {0, 147, 294, 441, 580},
    xticklabels = {$0$, $r_s$, $2r_s$, $3r_s$ },
    ytick = {-1, 0, 1},
    grid = both,
    grid style      = {line width=0.2pt, color=gray!25},
    major grid style= {line width=0.4pt, color=gray!35},
    clip = true,
    legend style = {
        at={(0.98,0.98)}, anchor=north east,
        font=\scriptsize, fill=white, fill opacity=0.85,
        draw=gray!50, inner sep=3pt, row sep=0pt,
        cells={anchor=west}
    },
]
\draw[gray!50, dashed, thin] (axis cs:0,0) -- (axis cs:580,0);
\draw[recline, dashed, line width=0.6pt, opacity=0.6]
    (axis cs:147,-1.5) -- (axis cs:147,1.5);
\draw[recline, dashed, line width=0.6pt, opacity=0.6]
    (axis cs:294,-1.5) -- (axis cs:294,1.5);
\draw[recline, dashed, line width=0.6pt, opacity=0.6]
    (axis cs:441,-1.5) -- (axis cs:441,1.5);
\addplot[modeone, line width=1.2pt, opacity=0.7, domain=0:580, samples=300]
    { -cos(deg(pi * x / 147)) };
\addlegendentry{$k_1$ (1-й пік)}
\addplot[modetwo, line width=1.2pt, opacity=0.7, domain=0:580, samples=300]
    { cos(deg(2*pi * x / 147)) };
\addlegendentry{$k_2$ (2-й пік)}
\addplot[modethree, line width=1.2pt, opacity=0.7, domain=0:580, samples=300]
    { -cos(deg(3*pi * x / 147)) };
\addlegendentry{$k_3$ (3-й пік)}
\addplot[modesum, line width=2pt, domain=0:580, samples=400]
    { ( -cos(deg(pi*x/147))
        + 0.6*cos(deg(2*pi*x/147))
        - 0.35*cos(deg(3*pi*x/147)) ) / 1.3 };
\addlegendentry{$\sum \delta_\gamma$ — спостережуваний розподіл}
\node[recline, font=\scriptsize, anchor=north] at (axis cs:147,-1.52)
    {$147\,\text{Мпк}$};
\end{axis}
\end{tikzpicture}
\caption{Просторовий «знімок» розподілу $\delta_\gamma(r)$ у момент
рекомбінації~--- суперпозиція мод, заморожених у різних амплітудах
(пор.\ з рис.~\ref{fig:acoustic_modes_A}).
Бурштинова крива~--- спостережуваний сигнал, що кодується у кутовий
спектр $\mathcal{D}_\ell^{TT}$.
Пунктирні вертикалі~--- масштаби $r_s$, $2r_s$, $3r_s \approx 147$~Мпк,
що відповідають акустичним пікам.}
\label{fig:acoustic_modes_B}
\end{figure}

Зручно увести скорочене позначення для середньої
(по напрямках) температурної флуктуації\footnote{%
Строго кажучи, $\Theta_0$~--- нульовий момент (монополь)
розкладу $\delta T/T$ за сферичними гармоніками по напрямку
імпульсу фотона; це стане зрозумілішим у розділі про
анізотропії CMB, де розглядається повна ієрархія
$\Theta_0, \Theta_1, \Theta_2,\ldots$}:
\begin{equation}\label{eq:Theta0_def}
  \Theta_0 \equiv \frac{\delta_\gamma}{4} = \frac{\delta T}{T},
\end{equation}
де останній зв'язок випливає із $\rho_\gamma \propto T^4$\footnote{%
$\delta\rho_\gamma = 4(\rho_\gamma/T)\,\delta T$,
тому $\delta_\gamma \equiv \delta\rho_\gamma/\bar\rho_\gamma
= 4\,\delta T/T$.
Множник~$4$~--- просто показник степеня.}.

Підставимо $\delta_\gamma = 4\Theta_0$ у рівняння
осцилятора~\eqref{eq:acoustic_preview}:
\begin{equation*}
  \frac{d}{d\eta}\bigl[(1+R)\cdot 4\Theta_0'\bigr]
  + \frac{k^2}{3}\cdot 4\Theta_0
  = -\frac{k^2}{3}(1+R)\,\Phi.
\end{equation*}
Множник $4$ скорочується. Розкриємо похідну в лівій частині:
\begin{equation*}
  (1+R)\,\Theta_0'' + R'\,\Theta_0'
  + \frac{k^2}{3(1+R)}\,\Theta_0
  = -\frac{k^2}{3}\,\Phi.
\end{equation*}
Ділимо на $(1+R)$ і отримуємо \textbf{рівняння гармонічного осцилятора}
із зовнішньою силою та згасаючим членом:
\begin{equation}\label{eq:Theta0_osc}
  \Theta_0'' + \frac{R'}{1+R}\,\Theta_0'
  + c_s^2 k^2\,\Theta_0 = -\frac{k^2}{3}\,\Phi,
  \qquad c_s^2 = \frac{1}{3(1+R)}.
\end{equation}
Тут $c_s$~--- швидкість звуку~\eqref{eq:sound_speed}.
Права частина~--- постійна зовнішня сила (зміщення рівноваги
гравітаційним потенціалом темної матерії).

Загальний розв'язок~--- косинус плюс зміщення рівноваги:
\[
  \Theta_0(k,\eta)
  = \bigl[\Theta_0(0) + (1+R)\Phi\bigr]\cos(k c_s \eta)
    - (1+R)\Phi.
\]
Для адіабатичних початкових умов ($\Theta_0(0) = -\Phi/2$,
$\Theta_0'(0) = 0$) у наближенні $R \ll 1$:
\begin{equation}\label{eq:Theta0_solution}
  \Theta_0(k,\eta) \approx \Theta_0(0)\cos(k c_s \eta).
\end{equation}
Це і є ті косинусні коливання, що зображені на
рис.~\ref{fig:acoustic_modes_A}: кожна мода коливається
зі своєю частотою $\omega_k = k c_s$, але всі стартують
з однакової фази~--- звідси когерентність і чіткі піки.

\begin{keybox}{Чому косинус, а не синус?}
Початкова умова від інфляції: $\delta_\gamma'(0) \approx 0$~---
плазма стартує зі стану максимального стиснення, без
початкової швидкості. Це як відпустити пружину зі стиснутого
стану~--- вона починає з максимуму, а не з нуля.
Якби початкова умова була $\delta_\gamma(0) = 0$, $\delta_\gamma'(0) \neq 0$~---
отримали б синус, і піки CMB стояли б в інших місцях.
Саме тому когерентність (однакова фаза для всіх $k$)~---
це пряма ознака інфляційного походження збурень: некогерентні
джерела (топологічні дефекти) дають випадкові фази і
розмазаний горб замість чітких піків.
\end{keybox}

%% ======================================================================
%% Рисунок А: еволюція мод у конформному часі
%% ======================================================================
\begin{figure}[h!]
\centering
\definecolor{modeone}{RGB}{232,89,60}
\definecolor{modetwo}{RGB}{29,158,117}
\definecolor{modethree}{RGB}{127,119,221}
\definecolor{recline}{RGB}{31,119,180}

\begin{tikzpicture}
\begin{axis}[
    width  = \linewidth,
    height = 0.42\linewidth,
    xmin = 0, xmax = 1.15,
    ymin = -1.3, ymax = 1.3,
    xlabel = {Конформний час $\eta$ (у.о.)},
    ylabel = {$\Theta_0(\eta)$},
    label style     = {font=\small},
    tick label style= {font=\footnotesize},
    axis lines = left,
    xtick = {0, 0.25, 0.5, 0.75, 1.0},
    xticklabels = {$0$, $0{,}25$, $0{,}5$, $0{,}75$, $\eta_{\rm rec}$},
    ytick = {-1, 0, 1},
    yticklabels = {$-1$, $0$, $1$},
    grid = both,
    grid style      = {line width=0.2pt, color=gray!25},
    major grid style= {line width=0.4pt, color=gray!35},
    clip = true,
    legend style = {
        at={(0.02,0.98)}, anchor=north west,
        font=\scriptsize, fill=white, fill opacity=0.85,
        draw=gray!50, inner sep=3pt, row sep=0pt,
        cells={anchor=west}
    },
]
\draw[gray!50, dashed, thin] (axis cs:0,0) -- (axis cs:1.15,0);
\draw[recline, dashed, line width=0.8pt]
    (axis cs:1.0,-1.3) -- (axis cs:1.0,1.3);
\addplot[modeone, line width=1.6pt, domain=0:1.15, samples=200]
    { cos(deg(pi * x)) };
\addlegendentry{$k_1$: $\lambda = 2r_s$, 1-й пік}
\addplot[modetwo, line width=1.6pt, domain=0:1.15, samples=200]
    { cos(deg(2*pi * x)) };
\addlegendentry{$k_2$: $\lambda = r_s$, 2-й пік}
\addplot[modethree, line width=1.6pt, domain=0:1.15, samples=200]
    { cos(deg(3*pi * x)) };
\addlegendentry{$k_3$: $\lambda = 2r_s/3$, 3-й пік}
\addplot[modeone,   only marks, mark=*, mark size=2.5pt]
    coordinates {(1.0, {cos(deg(pi))})};
\addplot[modetwo,   only marks, mark=*, mark size=2.5pt]
    coordinates {(1.0, {cos(deg(2*pi))})};
\addplot[modethree, only marks, mark=*, mark size=2.5pt]
    coordinates {(1.0, {cos(deg(3*pi))})};
\node[modeone,   font=\scriptsize, anchor=west] at (axis cs:1.02,-1.0) {1-й пік};
\node[modetwo,   font=\scriptsize, anchor=west] at (axis cs:1.02, 1.0) {2-й пік};
\node[modethree, font=\scriptsize, anchor=west] at (axis cs:1.02,-0.75){3-й пік};
\node[font=\scriptsize, recline, anchor=north west]
    at (axis cs:1.01, 1.28) {$\eta_{\rm rec}$};
\end{axis}
\end{tikzpicture}
\caption{Еволюція трьох характерних мод $\Theta_0(\eta) = \delta_\gamma/4$
у конформному часі (нормовано на $\eta_{\rm rec}$). Всі моди стартують з однакової фази~--- це наслідок інфляційного збурення $\mathcal{R}_k$ (§\,\ref{sec:quantum}).
Точки показують амплітуду кожної моди в момент «стоп-кадру»~---
рекомбінації $\eta_{\rm rec}$.}
\label{fig:acoustic_modes_A}
\end{figure}


%% --------------------------------------------------------
\section{Рекомбінація, звуковий горизонт і BAO}
%% --------------------------------------------------------

Акустичні коливання тривають до \emph{рекомбінації}
($z_\mathrm{rec} \approx 1100$, $T \approx 3000$~К):
протони захоплюють електрони, іонізація припиняється,
томсонівське розсіяння вимикається і фотони починають
вільно поширюватися. Плазма «заморожується» в тому стані,
в якому її застала рекомбінація.

Максимальна відстань, яку встигла пройти звукова хвиля
від початку гарячої епохи до рекомбінації,~---
\emph{радіус звукового горизонту}:
\begin{equation}\label{eq:sound_horizon}
  r_s(\eta_\mathrm{rec})
    = \int_{\eta_\mathrm{reh}}^{\eta_\mathrm{rec}} c_s(\eta)\,d\eta
    \approx 147\,\mathrm{Mpc},
\end{equation}
де $\eta_\mathrm{reh}$~--- момент реогріву (початок гарячої
епохи); оскільки $c_s \to 1/\sqrt{3}$ при ранніх часах,
нижня межа практично не впливає на результат і часто
пишуть $0$.
Це єдиний лінійний масштаб у задачі. Моди, що в момент
рекомбінації перебували у точці максимального стиснення
або розширення, проявляються у температурному спектрі
CMB як серія \emph{акустичних піків}~--- докладніше
у §\,\ref{sec:sachs_wolfe}.
Просторовий розподіл $\delta_\gamma(r)$ у цей момент~---
«стоп-кадр» заморожених коливань~--- показано на
рис.~\ref{fig:acoustic_modes_B}.

Але акустичні хвилі залишають слід не лише у фотонах.
Після рекомбінації фотони пішли геть, а баріонна речовина
«замерзла» у вигляді сферичної оболонки радіусом $r_s$
навколо кожної початкової неоднорідності темної матерії
(рис.~\ref{fig:bao_profile}). Це~--- \textbf{баріонні
акустичні осциляції (BAO)} \cite{eisenstein_2005}.
На рисунку показано профіль густини баріонів навколо
точкового первинного збурення: вісь~$r$~--- відстань
від центру збурення, крива~--- надлишок густини баріонів.
Чіткий пік на $r \approx r_s \approx 147$~Мпк~---
це «відлуння» звукової хвилі, замороженої у момент
рекомбінації.

\begin{figure}[ht!]
\centering
\begin{tikzpicture}
  \begin{axis}[
    width=\linewidth,
    height=0.7\linewidth,
    xmin=0, xmax=250,
    ymin=-0.05, ymax=1.1,
    xlabel={Відстань від первинного збурення, $r$ [Мпк]},
    ylabel={Контраст густини баріонів,
            $\Delta\rho_b/\bar{\rho}_b$},
    grid=both,
    grid style={line width=.1pt, color=gray!20},
    major grid style={line width=.2pt, color=gray!40},
    tick label style={font=\footnotesize},
    label style={font=\small},
    legend style={at={(0.95,0.95)}, anchor=north east,
                  font=\footnotesize, cells={anchor=west}},
    axis lines=left,
    clip=false
  ]
    \addplot[
      domain=0:240, samples=200,
      color=accent, line width=1.5pt
    ] {exp(-0.015*x) + 0.35*exp(-((x-147)/12)^2)};
    \addlegendentry{Профіль густини речовини}

    \draw[dashed, color=theoryred, line width=1pt]
      (axis cs:147,-0.05) -- (axis cs:147,0.9);
    \node[anchor=south, color=theoryred, font=\footnotesize]
      at (axis cs:147,0.9) {$r_s \approx 147\text{ Мпк}$};

    \node[anchor=west, color=gray, font=\scriptsize]
      at (axis cs:165,0.2) {Затухання Сілка};
    \draw[->, color=gray, >=stealth]
      (axis cs:162,0.18) -- (axis cs:155,0.1);
  \end{axis}
\end{tikzpicture}
\caption{Схематичний профіль поширення первинного
  адіабатичного збурення у баріонній компоненті плазми
  (модельна крива). Надлишок речовини на радіусі
  звукового горизонту $r_s$ відповідає сучасному
  масштабу BAO.}
\label{fig:bao_profile}
\end{figure}

У сучасну епоху масштаб BAO проявляється як фіксований
надлишок у просторовій кореляційній функції розподілу
галактик на відстані $\sim 147$~Мпк. Оскільки $r_s$
жорстко визначений лінійною фізикою раннього
Всесвіту~\eqref{eq:sound_horizon}, BAO~--- ідеальна
«стандартна лінійка»: вимірюючи цей масштаб на різних
червоних зміщеннях, можна реконструювати $H(z)$ і
параметри темної енергії незалежно від CMB.

\begin{keybox}{Чому BAO~--- «стандартна лінійка»?}
Уявіть що у стародавні часи хтось розставив мітки
кожні $147$~Мпк у просторі. Тепер, вимірюючи кутовий
розмір цих міток на різних відстанях ($z$), можна
дізнатися як простір між нами і ними розширювався~---
тобто $H(z)$. Головне: відстань $147$~Мпк визначена
фізикою раннього Всесвіту (швидкість звуку і час
рекомбінації), яка добре відома~--- тому «лінійка» надійна.
\end{keybox}

Той самий звуковий горизонт $r_s$, що визначає масштаб
BAO у розподілі галактик, відповідає і за положення
акустичних піків у температурному спектрі CMB. Як саме
ці піки формуються~--- і які інші ефекти на них
накладаються~--- розглядається у наступному розділі.

% ============================================================
% Задачі до розділу 6: Гідродинаміка плазми та акустичні осциляції
\section*{Задачі}
\addcontentsline{toc}{section}{Задачі}
% ============================================================

\begin{problem}{Рівняння неперервності та Ейлера для фотон-баріонної плазми}
Закон збереження тензора енергії-імпульсу $\nabla_\mu T^{\mu\nu}=0$ дає
рівняння неперервності \eqref{eq:continuity_general} та Ейлера \eqref{eq:euler_general}.

\begin{enumerate}
    \item Для фотонів ($w=1/3$, $c_s^2=1/3$) на субхабблових масштабах
          спростіть рівняння неперервності до вигляду $\delta_\gamma' = -\frac{4}{3}\theta + 4\Phi'$.
          Поясніть фізичний зміст кожного члена.
    \item Для фотон-баріонної рідини з параметром $R = 3\bar\rho_b/4\bar\rho_\gamma$
          рівняння Ейлера набуває вигляду \eqref{eq:euler_photon_baryon}.
          Поясніть, чому присутність баріонів ($R>0$) зменшує прискорення
          рідини під дією градієнта тиску.
    \item У границі $R \to 0$ (чиста фотонна плазма) покажіть, що з цих рівнянь
          випливає $\Theta_0'' + \frac{k^2}{3}\Theta_0 = -\frac{k^2}{3}\Phi$.
\end{enumerate}
\end{problem}

% -------------------------------------------------------------------

\begin{problem}{Акустичний осцилятор та початкові умови}
Рівняння для температурної флуктуації $\Theta_0 = \delta_\gamma/4$ у наближенні
$R \ll 1$ має вигляд $\Theta_0'' + c_s^2 k^2 \Theta_0 = -\frac{k^2}{3}\Phi$,
де $c_s^2 = 1/3$.

\begin{enumerate}
    \item Знайдіть загальний розв'язок цього рівняння при сталому $\Phi$.
    \item Інфляційні збурення задають початкові умови: $\Theta_0(0) = -\Phi/2$,
          $\Theta_0'(0) = 0$. Покажіть, що тоді $\Theta_0(\eta) = \Theta_0(0)\cos(k c_s \eta)$.
    \item Чому початкова фаза косинусна, а не синусна? Що б змінилося в кутовому
          спектрі CMB, якби $\Theta_0'(0) \neq 0$?
\end{enumerate}
\end{problem}

% -------------------------------------------------------------------

\begin{problem}{Звуковий горизонт та його значення}
Звуковий горизонт визначається як $r_s(\eta_{\text{rec}}) = \int_0^{\eta_{\text{rec}}} c_s(\eta)\,d\eta$.

\begin{enumerate}
    \item У радіаційно-доміновану епоху $c_s = 1/\sqrt{3}$ (в одиницях $c=1$),
          а конформний час $\eta \propto a$. Покажіть, що $r_s(\eta_{\text{rec}}) \approx \frac{1}{\sqrt{3}}\eta_{\text{rec}}$.
    \item Використовуючи $\eta = \int da/(a^2 H)$ та
          $H(a) = H_0 \sqrt{\Omega_r a^{-4} + \Omega_m a^{-3}}$,
          виконайте чисельне інтегрування або скористайтеся
          відомим результатом: $r_s \approx 147$ Мпк.
   \item Чому $r_s \approx 147$ Мпк є «стандартною лінійкою» для BAO?
          Як зміна $r_s$ вплинула б на положення акустичних піків у спектрі CMB?
\end{enumerate}
\end{problem}

% -------------------------------------------------------------------

\begin{problem}{Ефект баріонного завантаження}
Параметр $R = 3\bar\rho_b/4\bar\rho_\gamma$ зростає з часом, оскільки $\rho_b \propto a^{-3}$,
а $\rho_\gamma \propto a^{-4}$. На момент рекомбінації $R \sim 0.6$.

\begin{enumerate}
    \item Використовуючи точний розв'язок рівняння \eqref{eq:Theta0_osc},
          покажіть, що повне збурення, яке визначає температурну анізотропію,
          має вигляд $(\Theta_0 + \Phi)_{\text{rec}} = \Theta_0(0)\cos(k r_s) - R\Phi$.
    \item Поясніть, чому непарні акустичні піки ($\cos(k r_s) = -1$)
          підсилюються відносно парних ($\cos(k r_s) = +1$).
    \item Як за відношенням висот перших двох піків можна визначити
          баріонну густину $\Omega_b h^2$?
\end{enumerate}
\end{problem}

% -------------------------------------------------------------------

\begin{problem}{Чи можна побачити акустичні осциляції в розподілі галактик сьогодні?}

Після рекомбінації баріонна речовина «замерзає» на радіусі звукового
горизонту $r_s$ навколо кожної первинної неоднорідності темної матерії.
Це створює надлишок ймовірності знайти галактики на відстані $\sim r_s$
одна від одної.

\begin{enumerate}
    \item Оцініть фізичний масштаб $r_s$ у сучасну епоху (з урахуванням розширення).
          Виразіть його в Мпк.
    \item Який кутовий масштаб $\theta$ на небі відповідає цій відстані
          на червоному зсуві $z \sim 0.5$? ($d_A(z)$ — кутова відстань.)
    \item Чому вимірювання BAO на різних $z$ дозволяє визначити $H(z)$
          незалежно від калібрування відстаней?
\end{enumerate}
\end{problem}